%%======================================================================
\section{Challenges in Clinical Deployment}
\label{sec:challenges}
%%======================================================================

The TRL~3--4 state documented in Section~\ref{sec:applications} reflects three
categories of barriers. Technical gaps include manipulation precision, contact
safety during bipedal locomotion, and software validation against clinical
performance standards. Institutional gaps include regulatory pathway ambiguity,
liability framework absence, and actuarial non-coverage. A third category is
circular: regulatory approval requires clinical evidence; clinical evidence
requires Institutional Review Board (IRB)-sanctioned trials; IRB approval
requires safety assurance; and safety assurance requires performance standards
that do not yet exist.

\subsection{Manipulation and Dexterity for Clinical Tools}
\label{sec:manipulation}

Atar et al.~\cite{atar2025humanoids} is the only study to systematically document
manipulation performance gaps in a clinical task context. Atar et al.\ evaluated
a Unitree G1 with Inspire Gen4 dexterous hands across seven clinical procedures,
achieving approximately 70\% aggregate procedural success with non-clinician
teleoperators. The paper identifies two specific failure modes: force output
insufficiency (the actuator configuration cannot reliably generate the controlled
forces required for needle penetration or sustained grip against resistance), and
sensor sensitivity limitations (the tactile feedback loop is insufficient for
sub-Newton force resolution, preventing detection of the haptic signatures that
guide needle procedures in clinical practice).

These findings generalize across the procedure literature. Liang et
al.~\cite{liang2025lapsurgie} demonstrate that laparoscopic surgery requires
enforcing a remote-center-of-motion (RCM) constraint at the trocar site, an
engineering workaround for a capability that commercial humanoid arms do not
provide natively. Cho et al.~\cite{cho2026humanoid} identify three engineering
targets for endoscopic clinical translation: endoscope stability in confined
anatomical cavities, force compliance during tissue contact, and positioning
accuracy for reliable camera view.

Synthesizing across Atar et al., Liang et al., and Cho et al., clinical manipulation requires three
capabilities that current commercial humanoid arms do not yet reliably provide:
(1)~force precision at sub-Newton resolution for needle and catheter
procedures~\cite{atar2025humanoids}; (2)~tool-specific kinematic adaptation for
laparoscopes, rigid endoscopes, and ultrasound probes; and (3)~fingertip-level
tactile sensing for grip force discrimination, instrument slip detection, and
tissue compliance sensing. Operationally, progress requires force resolution at
sub-Newton levels for needle insertion procedures~\cite{atar2025humanoids},
fingertip sensor density on the order of $10^3$ sensing elements, a resolution
already reached by camera-based tactile skins such as
GelSight~\cite{yuan2017gelsight}, and teleoperation
round-trip latency at levels consistent with real-time force feedback
requirements. Surgical teleoperation research establishes that network latency
above 200~ms degrades task completion time and trajectory accuracy in
dexterous procedures~\cite{lum2009latency,xu2014latency,motiwala2025telesurgery},
providing the clinical benchmark for humanoid teleoperation system design.
Foundational work on dexterous manipulation documents the multi-dimensional
sensor requirements for reliable tool handling~\cite{okamura2000dexterous}, and
tactile-sensing reviews and high-resolution fingertip sensors define what
clinical-grade contact perception would
require~\cite{kappassov2015tactile,yuan2017gelsight}.

\subsection{Safe Physical Contact and Bipedal Gait Dynamics}
\label{sec:safety}

Making physical robot-human contact safe in human environments has been an
active research problem for two decades~\cite{alami2006safe}, and concrete
progress for humanoids has recently been demonstrated. Lu et
al.~\cite{lu2025gentlehumanoid} introduce a spring-based compliance framework for
the Unitree G1 that quantifiably reduces peak contact forces during quasi-static
tasks including sit-to-stand assistance and upper-body physical contact. This
approach builds on the series elastic actuator design~\cite{pratt1995elastic},
in which mechanical spring compliance decouples rapid force transients from the
actuated structure. This is the most technically grounded progress toward
patient-contact safety in the current literature. Hierarchical whole-body control
architectures~\cite{sentis2005synthesis} provide the control-theoretic foundation
for multi-priority task execution that safety-critical clinical interaction will
require.

The design target for patient-contact force limits is the Power and Force
Limiting (PFL) regime of ISO/TS~15066:2016~\cite{isots15066}, which provides
biomechanically derived thresholds for different body regions
(Table~A.2): for the face, quasi-static contact force is limited to 65~N and
pressure to 110~N/cm$^2$; for the chest, to 140~N and 120~N/cm$^2$
(sternum)~\cite{isots15066}. The GentleHumanoid framework demonstrates a
Unitree G1 approaching these thresholds for quasi-static upper-body contact.

The critical gap, however, is structural. ISO/TS~15066 was designed for
fixed-base collaborative robots (cobots whose base is anchored, whose motion is
constrained to a defined workspace, and whose approach velocities are programmed).
It does not model contact forces arising from bipedal locomotion. A bipedal
humanoid navigating a clinical environment introduces four risk scenarios that
ISO/TS~15066 does not bound: (1)~balance corrections, in which whole-body torque
adjustments generate rapid impulse forces that can contact nearby persons;
(2)~stumble recovery, in which momentum transfer creates swinging limbs and torso
rotation at velocities above the standard's contact speed assumptions;
(3)~narrow-corridor navigation, in which swinging arms during gait will contact
bed rails, IV poles, and patient limbs in spaces with less than one body-width of
clearance; and (4)~postural transitions, in which a full-size humanoid (35~kg or
more) rising to or lowering from a working height near a supine patient generates
dynamics that quasi-static force limits do not characterize.

None of the reviewed clinical studies addresses this category of risk, and neither
ISO/TS~15066~\cite{isots15066}, ISO~13482~\cite{iso13482}, nor
ISO~10218-1/2~\cite{iso10218_1,iso10218_2} was designed to model or bound
gait-induced contact events. This is an open safety engineering problem.
Resolution requires either new standards specific to bipedal humanoid robots in
patient-contact environments, or architectural solutions (proximity-triggered
gait interrupts, patient-detection zones constraining locomotion near occupied
beds, whole-body collision prediction for gait trajectories) that have not yet
been developed in the published literature.

\subsection{Human-Robot Trust and Institutional Adoption}
\label{sec:trust}

Trust and adoption operate at three distinct layers. The first layer is patient
and clinician acceptance. Technology acceptance frameworks established for
information systems~\cite{venkatesh2003user} and their extensions for assistive
social robots in elderly care~\cite{heerink2010assessing,broadbent2017interactions}
provide the theoretical basis for understanding humanoid adoption barriers.
Studies of robot acceptability in domestic and care settings further establish
that transparency, predictability, and social role clarity determine
acceptance~\cite{young2009toward}.
Trust in human-robot interaction is a multidimensional construct influenced by
robot performance, human factors, and environmental
context~\cite{hancock2011meta,hancock2011trust,sciutti2018humanizing}. The healthcare-specific question of whether
humanoid form factors in clinical settings are perceived as appropriate has been
characterized as a strategic and ethical
question~\cite{ozturkcan2022humanoid}. Robinson et al.~\cite{robinson2023brief}
show 53\% recruitment uptake for a humanoid wellbeing intervention in a general
population (pilot RCT, n\,=\,230). Sayis and Gunes~\cite{sayis2024technology}
show the humanoid condition was the only condition producing measurable mood
improvements in a therapeutic writing task (n\,=\,42). Both results are bounded:
the clinical populations central to humanoid deployment differ substantially from
general and educational populations in these studies.

The second layer is co-design and workflow integration. Ghosh et
al.~\cite{ghosh2024usercentered} demonstrate that iterative co-design with care
staff improves system usability, a methodological contribution that does not
substitute for clinical evidence of safety and effectiveness.

The third and most institutionally consequential layer is the liability gap. No
current healthcare liability framework (US tort law, UK clinical negligence, EU
product liability) cleanly assigns accountability for harm caused by an autonomous
robotic action in a care setting: as autonomy increases, the agency, control, and
foreseeability tests that negligence and product liability doctrines rely on break
down, leaving a recognized accountability
gap~\cite{eldakak2024civil,solaiman2024regulating}. This legal ambiguity deters adoption regardless
of technical capability: hospital risk departments cannot quantify institutional
exposure for systems operating in legally uncharted territory. Imtiaz and
Khan~\cite{imtiaz2024perceptions} document this precisely: 269 nursing home
administrators cite cost, reduction in human contact, and unproven clinical
effectiveness as primary barriers.

\subsection{Regulatory Pathway}
\label{sec:regulatory}

The regulatory challenge is not hostility to this technology class; it is a
jurisdictional gap that existing frameworks leave precisely where clinical
humanoids sit.

ISO~13482:2014~\cite{iso13482} (Safety Requirements for Personal Care Robots) is
the most directly applicable safety standard; it governs robots providing
physical assistance in human environments. However, \S{}1.4 of ISO~13482
explicitly excludes devices regulated as medical devices. A humanoid robot performing
clinical tasks almost certainly meets the medical device definition under the US
Federal Food, Drug, and Cosmetic Act \S{}201(h)~\cite{fdcact201h} and EU MDR
Article~2(1)~\cite{eumdr2017}. This creates a hard gap: the system
is likely too clinical for ISO~13482 but no ISO medical robotics standard covers
general-purpose bipedal humanoid assistants.

In the United States, the De~Novo classification pathway is the only viable
regulatory entry route. Because no legally marketed device of comparable
intended use and technology exists, there is no predicate to support a 510(k)
substantial-equivalence claim, which is precisely the case the De~Novo route was
created to address~\cite{fda2021denovo}. A De~Novo submission requires
the applicant to propose Special Controls from scratch: accepted performance
standards for fall prevention in patient-care environments (an area where the
patient-safety literature is well developed for human falls but silent on
robot-induced ones)~\cite{oliver2007strategies,cameron2012interventions}, patient
contact force limits for dynamic gait scenarios, teleoperation latency bounds for specific
procedure types, and AI/ML software validation methodology consistent with FDA's
AI/ML SaMD action plan~\cite{fda2021aiml} and December~2024 PCCP
guidance~\cite{fda2024pccp}. None of these currently exist as accepted standards. This is not procedurally
impossible: De~Novo has established new device categories for surgical planning
software and exoskeletal rehabilitation devices, but it represents a development
and regulatory investment no humanoid manufacturer has yet made.

The structural problem is a circular dependency: approval requires clinical
evidence; human-subject evidence requires IRB approval; IRB approval requires
safety assurance; safety assurance requires performance standards; performance
standards require regulatory guidance. The productive path through this circle is
a staged evidence generation sequence (in vitro feasibility, cadaver study,
IRB-approved human-subject feasibility trial) producing the data needed for a
De~Novo submission. Cho et al.~\cite{cho2026humanoid} is the current frontier of
that sequence.

In the EU, a humanoid clinical assistant would likely be classified as Class~IIa
under EU MDR Annex~VIII, with autonomous surgical assistance escalating to
Class~IIb or~III~\cite{eumdr2017}. EU AI Act Regulation 2024/1689 additionally classifies any
AI-driven clinical humanoid as a high-risk AI system~\cite{euaiact2024,solaiman2024regulating}. Under Article~113,
high-risk AI obligations apply from 2~August~2026, while AI embedded in products
regulated under existing Union law such as the MDR transitions on
2~August~2027~\cite{euaiact2024}. Any EU clinical humanoid market entry must satisfy both MDR and AI Act
simultaneously, a dual compliance burden for which no completed template yet
exists.

Table~\ref{tab:regulatory} compares the four principal frameworks across the key
dimensions that determine market entry strategy.

\begin{table*}[t]
  \captionsetup{justification=centering,skip=2pt}
  \caption{Regulatory Framework Comparison for Clinical Humanoid Robots (March 2026).}
  \label{tab:regulatory}
  \resizebox{\linewidth}{!}{%
  \begin{tabular}{@{}lllllp{4cm}@{}}
    \toprule
    \textbf{Framework} & \textbf{Jurisdiction} & \textbf{Classification} &
    \textbf{Pathway} & \textbf{Timeline (est.)} & \textbf{Key Requirements / Gap} \\
    \midrule
    FDA 510(k)      & US     & Class I/II  & Substantial equivalence to predicate & 3--12 months  & No predicate exists for clinical humanoid; inapplicable \\
    FDA De~Novo     & US     & Class II (new category) & Novel device; propose Special Controls & 12--24 months & Must define fall-prevention, PFL, latency, AI/ML Special Controls from scratch \\
    FDA PMA         & US     & Class III   & Clinical trial evidence required & 3--7 years & Overkill for initial classification; may apply for autonomous surgical \\
    EU MDR Class IIa & EU    & Medium risk & Notified Body QMS + clinical data & 18--36 months & Requires clinical investigation under Art.\ 62; PMCF post-market \\
    EU MDR Class IIb & EU   & Higher risk & Notified Body + full clinical eval.\ & 24--48 months & Autonomous surgical or patient-contact without operator oversight \\
    EU AI Act (High Risk) & EU & High-risk AI & Conformity assessment; Art.\ 6 & From Aug 2026 & Parallel obligation to MDR; dual compliance burden for clinical humanoid \\
    ISO 13482:2014  & Global & Personal care robot safety & Voluntary standard & N/A & \textbf{Explicitly excludes medical devices}; no bipedal gait dynamic PFL spec \\
    ISO/TS 15066:2016 & Global & Collaborative robot HRC & Voluntary; used in EU MDR & N/A & Designed for fixed-base collaborative arms; no dynamic gait scenario coverage \\
    \bottomrule
    \multicolumn{6}{@{}l@{}}{\textbf{PMS} = post-market surveillance; \textbf{PMCF} = post-market clinical follow-up; \textbf{SaMD} = Software as a Medical Device; \textbf{NB} = Notified Body.} \\
    \multicolumn{6}{@{}l@{}}{No framework provides humanoid-specific guidance; all entries represent the authors' interpretation of existing rules applied to this novel device class.} \\
  \end{tabular}%
  }
\end{table*}

\subsection{Software Infrastructure and AI Validation}
\label{sec:software}

The leading VLA models, including openpi/$\pi_0$ (10,551~stars), OpenVLA (5,479~stars),
and LeRobot (22,179~stars), represent the highest-capability general-purpose
robot learning infrastructure publicly available. For general-purpose humanoid
manipulation research, this infrastructure is genuinely strong.

What is absent is the layer between this general infrastructure and clinical
deployment. As identified in Section~\ref{sec:accessibility}, the four critical gaps
with no current open-source solution are a clinical safety layer, clinical
workflow integration, an open clinical demonstration dataset, and a validated
medical simulation environment.

The AI validation gap extends to deployment readiness. The leading VLA models
achieve high success rates on standard manipulation benchmarks, but those
benchmarks share no tasks with the clinical procedure requirements documented by
Atar et al.~\cite{atar2025humanoids}, and the shift from benchmark to clinic
involves not only visual and geometric differences but also contact physics, tool
compliance, and procedural sequencing fundamentally different from household
manipulation. A success rate acceptable for a benchmark grasping task is not
acceptable for needle insertion or catheter placement, where the approximately
70\% aggregate success observed on real clinical tasks~\cite{atar2025humanoids}
translates directly into patient risk. This mirrors a broader pattern in medical
AI, where models that perform well in retrospective evaluation degrade under
prospective deployment and distribution
shift~\cite{topol2019high,rajpurkar2022ai,coiera2019last,sendak2020real}.

Beyond raw task performance, learned control policies raise a verification and
validation problem that current infrastructure does not solve. A VLA policy is a
function defined over a very large parameter space, and its safety envelope cannot
be exhaustively characterized by testing; yet regulatory clearance through the
De~Novo route would require exactly such evidence, expressed as accepted Special
Controls and, for models that keep learning after deployment, a predetermined
change control plan~\cite{fda2021aiml,fda2024pccp}. No humanoid software stack
today supplies the artifacts this demands: a documented safety case, a bounded and
continuously monitored operational design domain, and validated behavior under the
sensor noise, occlusion, and contact uncertainty of a real ward. The open
demonstration datasets and validated medical simulation environments that would
let such a safety case be assembled do not yet exist for clinical humanoid tasks,
as identified in Section~\ref{sec:accessibility}.

Security compounds the problem. No humanoid software stack has addressed
IEC~62443 Security Level~2 requirements, the minimum threshold for networked
systems in hospital environments; a robot connected to hospital networks while in
physical contact with patients presents a combined digital and physical attack
surface for which no security validation specific to humanoids has been
published.

Table~\ref{tab:techgaps} synthesizes the five principal technical and
translational barriers identified in this section.

\begin{table*}[t]
  \captionsetup{justification=centering,skip=2pt}
  \caption{Technical and Translational Gap Summary (March 2026).}
  \label{tab:techgaps}
  \footnotesize
  \resizebox{\linewidth}{!}{%
  \begin{tabular}{@{}>{\raggedright\arraybackslash}p{2.60cm}
                  >{\raggedright\arraybackslash}p{4.70cm}
                  >{\raggedright\arraybackslash}p{4.80cm}
                  >{\raggedright\arraybackslash}p{2.70cm}
                  >{\raggedright\arraybackslash}p{3.15cm}@{}}
    \toprule
    \textbf{Capability Gap} & \textbf{Current State} & \textbf{Required for TRL~5} &
    \textbf{Blocking Domains} & \textbf{Key Evidence / Standard} \\
    \midrule
    Force precision \& tactile sensing &
      Insufficient sub-N resolution; no fingertip tactile array on commercial platforms &
      Sub-Newton force feedback; continuous fingertip sensing at clinical manipulation precision &
      Clinical procedures, Rehabilitation &
      \cite{atar2025humanoids}; \cite{lum2009latency} \\[4pt]
    Gait contact safety &
      ISO/TS~15066 quasi-static thresholds approached for upper body only~\cite{lu2025gentlehumanoid} &
      Instrumented characterization of gait-induced force profiles; accepted standard for bipedal PFL &
      All physical-contact domains &
      ISO/TS~15066~\cite{isots15066}; ISO~13482~\cite{iso13482} \\[4pt]
    Tool-specific kinematics &
      No native RCM; no laparoscope-specific axis on commercial arms &
      RCM enforcement or software analog; tool-path control within trocar constraints &
      Clinical procedures &
      \cite{liang2025lapsurgie}; \cite{cho2026humanoid} \\[4pt]
    VLA clinical generalization &
      VLA benchmarks: household manipulation only; no clinical demonstration dataset exists &
      Open clinical demonstration dataset; VLA fine-tuned on clinical task protocols with validated performance &
      Clinical procedures, all autonomous roles &
      \cite{openxembodiment2023}; \cite{zhao2023act} \\[4pt]
    Clinical safety standard \& regulatory predicate &
      No accepted standard for bipedal gait in patient-contact environments; no De~Novo predicate &
      ISO scope revision; IRB feasibility study; FDA De~Novo with proposed Special Controls &
      All domains &
      \cite{iso13482}; \cite{fda2021aiml}; \cite{fda2024pccp} \\
    \bottomrule
    \multicolumn{5}{@{}p{17.95cm}@{}}{Each row identifies a barrier from Section~\ref{sec:challenges}, its current state, the threshold required for TRL~5 advancement, blocked clinical domains, and the primary evidence or standard defining it.} \\
    \multicolumn{5}{@{}p{17.95cm}@{}}{No single gap is insurmountable; the barriers are sequentially dependent in the order listed.} \\
  \end{tabular}%
  }
\end{table*}
